\scanpage{14}
\mspage{-0.19-}
\[
\phi(\alpha X+\beta Y,F_4(X,Y))=F_4(-\beta,\alpha).
\]
Then it is easy to check that $(\mathbb A^1,\phi)$ is a geometric quotient of $X$ by $\mathrm{SL}(2)$. Finally, the morphism $\unsure{\overline{\underline{\psi}}|\overline{\underline{\Psi}}}$ is not even closed. Let $Z$ be the closed subscheme of $G\times X$ whose closed points are the pairs:
\[
\begin{pmatrix}0&-\lambda^{-1}\\\lambda&0\end{pmatrix}\times(\lambda X+Y,X^2Y^2)
\]
for $\lambda\in\mathbb C-(0)$. The image under $\unsure{\overline{\underline{\psi}}|\overline{\underline{\Psi}}}$ is the above point is:
\[
(-\lambda X+Y,X^2Y^2)\times(\lambda X+Y,X^2Y^2).
\]
In the closure of this set is the extra point:
\[
(Y,X^2Y^2)\times(Y,X^2Y^2).
\]
\uline{Lemma} \uline{0.5}: Let $\sigma$ be an action of $\mathbb G_m$ on an algebraic scheme $X$, all over a field $k$. Assume that $X$ admits an immersion in projective space $\mathbb P_n$ for which the action $\sigma$ extends to an action of $\mathbb G_m$ on $\mathbb P_n$. Then:
\[
\sigma\text{ is proper}\Longleftrightarrow\begin{cases}\text{(i)}&\sigma\text{ is separated},\\\text{(ii)}&S_1(X)=\phi.\end{cases}
\]
\uline{Proof}: The implication $\Longrightarrow$ is clear. Conversely, assume $\sigma$ is separated and $S_1(X)=\phi$. To prove that the morphism $\unsure{\overline{\underline{\psi}}|\overline{\underline{\Psi}}}$ is proper, we use the valuative criterion (EGA 2, \S7). Let $R$ be a valuation ring over $k$, and let $K$ be its quotient field. We may assume that $R$ contains an algebraically closed overfield of $k$, $\Omega$, which is isomorphic to its residue field (Remark 7.3.9, EGA 2). Now suppose $\unsure{\phi|\rho}\times\xi$ is a $K$-valued but not $R$-valued point
\mspage{13}
